\einheitenkopf{Einheit 4 von 5 · Kenngrößen}
\erg{29}{a) Spannweite \quad b) Modalwert \quad c) Median \quad d) Mittelwert \quad e) Modalwert}
\erg{30}{a) $3\ 4\ 6\ 7\ 9$ \quad b) $9\ 9\ 11\ 12\ 14\ 15$ \quad c) $1{,}8\ \ 2{,}0\ \ 2{,}5\ \ 3{,}1$ \quad d) $98\ 98\ 103\ 105\ 112$ \quad e) $0{,}7\ \ 0{,}8\ \ 0{,}9\ \ 1{,}2\ \ 1{,}5$}
\erg{31}{a) Min 2, Max 11 \quad b) Min 15, Max 30 \quad c) Min $2{,}9$, Max $4{,}1$ \quad d) Min $7\,^\circ$C, Max $14\,^\circ$C \quad e) $11-2=9$ \quad f) $9{,}2-8{,}5=0{,}7$ \quad g) 6 \quad h) geordnet $1\ 2\ 5\ 8\ 9$ → Median 5 \quad i) geordnet $11\ 14\ 17\ 20$ → $(14+17):2 = 15{,}5$ \quad j) Min 20\,mm, Max 60\,mm, Spannweite 40\,mm}
\erg{32}{a) 5 \quad b) 5 \quad c) 20 \quad d) 3 \quad e) $52 : 4 = 13$ \quad f) $51{,}4 : 6 = 8{,}5\overline{6} \approx 8{,}6$ \quad g) $4390 : 17 \approx 258{,}2$, also rund 258 Besucher \quad h) $344 : 8 = 43$ \quad i) $6 \cdot 15 = 90$; $90 - 71 = 19$ \quad j) $306 : 9 = 34$ Jahre; ohne die älteste (52) und die jüngste (16) Person: $238 : 7 = 34$ Jahre. Der Mittelwert ändert sich nicht, weil die beiden wegfallenden Werte zusammen $68 = 2 \cdot 34$ ergeben, also genau zweimal den bisherigen Mittelwert.}
\erg{33}{a) ja \quad b) nein, der Modalwert ist 4 \quad c) ja \quad d) nein, der Mittelwert ist $36 : 5 = 7{,}2$}
\erg{34}{a) Nina hat den mittleren Wert der unsortierten Liste genommen. Erst ordnen: $3\ 5\ 7\ 9\ 12$; der Median ist 7. \quad b) geordnet $11\ 14\ 18\ 20\ 22\ 25$ → $(18+20):2 = 19$}
\erg{35}{a) Der eine sehr hohe Verdienst zieht den Mittelwert nach oben, obwohl ihn niemand außer der Chefin erhält. Der Median hängt nur von der Reihenfolge ab und liegt deshalb bei 2400\,€. \quad b) Nein. Summe und Anzahl wachsen so, dass der Quotient gleich bleibt: Der neue Wert gleicht dem bisherigen Durchschnitt und verschiebt ihn deshalb nicht.}
